\chapter{Relación con la titulación}

\section{Competencias generales y habilidades}

De las competencias genéricas de referencia para las Prácticas Académicas Externas, identifico a continuación las que desarrollé con mayor intensidad durante el periodo de prácticas, junto con el ejemplo concreto que las ilustra:

\begin{itemize}
  \item \textbf{Iniciativa.} Busqué de forma autónoma, sin instrucción directa, el artículo científico que describe la metodología HAVEN, y validé sus principios mediante pruebas propias antes de proponerla formalmente como base metodológica del proyecto.
  \item \textbf{Innovación y creatividad.} Construí un prototipo funcional basado en HAVEN como propuesta propia, aprovechando la libertad que me concedió mi tutor de la entidad colaboradora para explorar herramientas y estrategias.
  \item \textbf{Planificación.} Coordiné mi sistema de generación de testbench (HAVEN) con el sistema de generación de test plan desarrollado por mi compañero de prácticas, manteniendo ambos sincronizados pese a los cambios constantes de criterios; gestioné mi propio tiempo de trabajo bajo metodología ágil, con reuniones diarias de 15 minutos.
  \item \textbf{Trabajo en equipo.} Colaboré de forma continua con mi compañero de prácticas y con el asesor técnico externo, buscando de forma constante el consenso sobre la estructura del test plan para evitar ambigüedad entre los dos sistemas que desarrollamos.
  \item \textbf{Habilidades comunicativas y comprensión interpersonal.} Me comuniqué habitualmente en inglés con el asesor técnico y con la documentación interna de la Fundación; expuse el prototipo HAVEN ante los equipos de diseño y verificación el 29 de junio.
  \item \textbf{Responsabilidad.} Cumplí mis entregas bajo un cronograma real de empresa, sensiblemente distinto al entorno académico, hasta la entrega final empaquetada al equipo de verificación de INNOVA IRV.
  \item \textbf{Autoconfianza.} Defendí técnicamente mi propuesta HAVEN ante los equipos de diseño y verificación, y ante la dirección de la Fundación en la presentación del estudio comparativo entre modelo local y modelo en la nube.
\end{itemize}

\section{Competencias específicas}

La conexión entre mis prácticas y las competencias específicas de la titulación fue, principalmente, la de la Inteligencia Artificial generativa aplicada a un problema de dominio específico (la verificación de hardware), enlazando con los contenidos finales de la asignatura de \emph{Aprendizaje Automático}. De forma adicional, los conocimientos que adquirí sobre la plataforma AWS durante la asignatura \emph{Bases para el Almacenamiento y Análisis del Dato} resultaron fundamentales para el uso de la API de Amazon Bedrock, que empleé para realizar las llamadas al modelo Claude Sonnet 4.6 en el desarrollo del proyecto.
